\begin{trapbox}
	\begin{description}[style=nextline, leftmargin=2.8em, labelsep=0.5em, itemsep=0.35em]
		\item[\textbf{\textcolor{CTrap}{易错点一（混淆射影角）}}] 误将直线与平面内不过垂足的直线所成角当作线面角。
		\item[\textbf{\textcolor{CTrap}{正确理解（射影定义）：}}] 线面角必须是直线与它在平面内的射影所成的锐角，射影必过垂足。
		\item[\textbf{\textcolor{CTrap}{错因分析（射影不清晰）：}}] 没有明确射影的概念，只凭直观选取平面内的任意直线。
	\end{description}

	\medskip
	\begin{description}[style=nextline, leftmargin=2.8em, labelsep=0.5em, itemsep=0.35em]
		\item[\textbf{\textcolor{CTrap}{易错点二（公式用错）}}] 误认为 $\cos\theta = \dfrac{QH}{PQ}$，将正弦与余弦混淆。
		\item[\textbf{\textcolor{CTrap}{正确理解（正弦公式）：}}] 线面角的正弦公式是 $\sin\theta = \dfrac{QH}{PQ}$，注意是正弦不是余弦。
		\item[\textbf{\textcolor{CTrap}{错因分析（三角函数混淆）：}}] 在立体几何中，线面角常与三角函数结合，但学生易混淆 $\sin$ 和 $\cos$。
	\end{description}

	\medskip
	\begin{description}[style=nextline, leftmargin=2.8em, labelsep=0.5em, itemsep=0.35em]
		\item[\textbf{\textcolor{CTrap}{易错点三（范围遗忘）}}] 认为线面角可以是钝角（如 $120^\circ$）。
		\item[\textbf{\textcolor{CTrap}{正确理解（范围限制）：}}] 线面角 $\theta$ 的范围是 $0^\circ \le \theta \le 90^\circ$，必为锐角或直角。
		\item[\textbf{\textcolor{CTrap}{错因分析（定义理解偏差）：}}] 误将直线与平面所成角当作直线与平面法线的夹角或其补角，忽略了锐角原则。
	\end{description}
\end{trapbox}
